\section{Expansions}
\label{sec:xml:expansions}
This section defines the polynomial expansions used on each of the defined
geometric composites and variables. Expansion entries specify the number of
modes and the expansion type, or a full list of data of basis type, number of
modes, points type and number of points. The short-hand version has the
following form

\begin{lstlisting}[style=XMLStyle]
<E COMPOSITE="C[0]" NUMMODES="5" FIELDS="u" TYPE="MODIFIED" />
\end{lstlisting}

or, if we have more then one variable we can apply the same basis to all using

\begin{lstlisting}[style=XMLStyle]
<E COMPOSITE="C[0]" NUMMODES="5" FIELDS="u,v,p" TYPE="MODIFIED" />
\end{lstlisting}

Each type chooses its own quadrature. \inltt{MODIFIED} and
\inltt{GLL\_LAGRANGE} take one point more than modes on
Gauss-Lobatto-Legendre points in the tensor-product directions and as many
points as modes on Gauss-Radau points in the collapsed ones.
\inltt{ORTHOGONAL} takes as many points as modes on the Gauss family:
Gauss-Legendre in the tensor-product directions and in a direction collapsed
once (the second of a triangle or tetrahedron, the third of a prism), and
Gauss-Jacobi(1,0) in a direction collapsed twice (the third of a tetrahedron
or pyramid). The orthogonal mass matrix is then integrated exactly on every
shape but the pyramid, every direction has the same count, and the element
and trace points coincide wherever the rule is Gauss-Legendre.

The short-hand form takes an optional \inltt{NUMPOINTS}, the number of
quadrature points in every direction, with the same expression syntax as
\inltt{NUMMODES}; it replaces the type's own count and keeps its point types:

\begin{lstlisting}[style=XMLStyle]
<E COMPOSITE="C[0]" NUMMODES="4" NUMPOINTS="5" FIELDS="u" TYPE="ORTHOGONAL" />
\end{lstlisting}

The expansion basis can also be specified in detail as a combination of
one-dimensional bases, and thus the user is able to, for example, increase the
quadrature order. For tet elements this takes the form:

\begin{lstlisting}[style=XMLStyle]
<E COMPOSITE="C[0]" 
   BASISTYPE="Modified_A,Modified_B,Modified_C" 
   NUMMODES="3,3,3"
   POINTSTYPE="GaussLobattoLegendre,GaussRadauMAlpha1Beta0,GaussRadauMAlpha2Beta0"
   NUMPOINTS="4,3,3"
   FIELDS="u" />
\end{lstlisting}

and for prism elements:

\begin{lstlisting}[style=XMLStyle]
<E COMPOSITE="C[1]" 
   BASISTYPE="Modified_A,Modified_A,Modified_B" 
   NUMMODES="3,3,3"
   POINTSTYPE="GaussLobattoLegendre,GaussLobattoLegendre,GaussRadauMAlpha1Beta0"
   NUMPOINTS="4,4,3"
   FIELDS="u" />
\end{lstlisting}

The expansions can be defined with a list of \inltt{<E>} elements (e.g., to
represent different polynomial orders for different variables or to address
different composites). The user can define a default expansion field by entering
\inltt{<E>} tags without the \inltt{FIELDS} attribute. The default expansion is
used to define any variables not explicitly listed in the \inltt{<E>} entries.
In the following example, the default expansion is used to define the expansions
for the composites C[0], C[1] and C[2]:

\begin{lstlisting}[style=XMLStyle]
<E COMPOSITE="C[0-2]" NUMMODES="5" TYPE="MODIFIED" />
<E COMPOSITE="C[3]"   NUMMODES="4" TYPE="MODIFIED" FIELDS="u,v"/>
<E COMPOSITE="C[3]"   NUMMODES="3" TYPE="MODIFIED" FIELDS="p"/>
\end{lstlisting}

The expansions of each field should be defined only once for each composite.

Where the mesh names its composites, a name may be given in place of an ID, so
that \inltt{COMPOSITE="C[fluid]"} selects the composite named \inltt{fluid};
see section~\ref{sec:xml:composites}.
